\documentclass{article}
\usepackage{graphicx}
\usepackage[margin=1.5cm]{geometry}
\usepackage{amsmath}
\usepackage{enumitem}

\setlist{nosep,leftmargin=*}

\begin{document}
\twocolumn

\title{Chapter 5 In-Class Exercise: Searching a Matrix}
\author{Prof. Jordan C. Hanson}
\date{}

\maketitle

\section{Matrices Using Lists}

\begin{enumerate}

\item A two-dimensional array of values can be represented in Python
using a \textbf{list of lists}.  Each inner list represents one row.  Begin with the matrix

\[
A =
\begin{bmatrix}
4 & 7 & 2 & 9 \\
1 & 5 & 8 & 3 \\
6 & 0 & 4 & 7
\end{bmatrix}.
\]

Represent this matrix in Python using

\begin{verbatim}
A = [
    [4, 7, 2, 9],
    [1, 5, 8, 3],
    [6, 0, 4, 7]
]
\end{verbatim}

\begin{itemize}

\item The first index selects a row and the second index selects a column.  Try

\begin{verbatim}
print(A[0][0])
print(A[1][2])
print(A[2][3])
\end{verbatim}

\item Remember that Python begins counting at zero.  Therefore, \texttt{A[1][2]} is the element in the second row and third column.

\end{itemize}

\end{enumerate}


\section{Nested Loops}

\begin{enumerate}

\item A loop inside another loop is called a \textbf{nested loop}.  The outer loop can move through the rows, while the inner loop moves through the columns. Use

\begin{verbatim}
for row in range(3):
    for column in range(4):
        print(A[row][column])
\end{verbatim}

\begin{itemize}

\item How many times does the outer loop execute?
\item How many times does the inner loop execute?
\item How many total matrix elements are printed?
\item Modify the print statement to display the row, column, and value:

\begin{verbatim}
print(row, column,
      A[row][column])
\end{verbatim}

\end{itemize}

\end{enumerate}


\section{Searching for a Value}

\begin{enumerate}

\item Now search the matrix for a particular value. Begin with

\begin{verbatim}
target = 7
\end{verbatim}

Use nested loops and an \texttt{if} statement:

\begin{verbatim}
for row in range(3):
    for column in range(4):

        if A[row][column] == target:
            print("Found at",
                  row, column)
\end{verbatim}

\begin{itemize}

\item Run the program.  How many times does the value \texttt{7} occur?
\item At which row and column indices does it occur?
\item Change

\begin{verbatim}
target = 7
\end{verbatim}

to several other values and repeat the search.

\end{itemize}

\end{enumerate}


\section{In-Class Exercise}

\begin{enumerate}

\item Create the matrix

\[
B =
\begin{bmatrix}
12 & 5 & 8 & 3 & 11 \\
7 & 14 & 2 & 9 & 6 \\
4 & 10 & 13 & 5 & 1 \\
8 & 3 & 6 & 12 & 7
\end{bmatrix}.
\]

Represent it as a list of four lists.

\begin{itemize}

\item Ask the user for an integer to search for:

\begin{verbatim}
target = int(input(
    "Value to search for: "))
\end{verbatim}

\item Use nested \texttt{for} loops to search every row and column.

\item Whenever the value is found, print

\begin{verbatim}
Value found at row __ column __
\end{verbatim}

using the actual row and column indices.

\item Create the Boolean variable

\begin{verbatim}
found = False
\end{verbatim}

before beginning the search. Whenever a match occurs, change it to

\begin{verbatim}
found = True
\end{verbatim}

\item After both loops have finished, use

\begin{verbatim}
if not found:
    print("Value not found")
\end{verbatim}

so that the program also correctly handles values that are not
inside the matrix.

\item Test your program using one value that occurs once, one value
that occurs more than once, and one value that does not occur.

\end{itemize}

\end{enumerate}

\end{document}